
Comme \(L_Y\) est \([-a_1]\)-régulier, \(h^2(L(-a_1-2))=0\) par 1.4; d'où \(h^1(L(-a_1-2))=h^0(L(-a_1-2))-\chi(L(-a_1-2))\). Or, \(h^0(L(-a_1-2))\le a_2\binom{-a_1}{2}\) par 1.8 (i) puisque \(L\) est un \((0,0,a_2)\)-faisceau par 2.5. Par suite
\[
0\le h^1(L(-a_1-2))\le -a_1(-a_1-1)-a_0.
\tag{2.7.2}
\]
Enfin, 2.7.2 donne (ii), et 2.7.1 et 2.7.2 donnent (i).

\textbf{Lemme 2.8.} \emph{On suppose que \(r=2\) et \(d(L)\le0\). Soient \(\delta=d(L)\) et \(\beta=\delta(e_2-1)+e_1\). Alors} :
\[
a_0\ge
\langle c_1(L)^2\rangle+2e_0+
\sum_{i=1}^{2}(2e_i-a_i)\binom{\delta-1+i}{i}
-\beta(\beta+1).
\]

En effet, soit \(\chi\bigl(L^{\otimes m}(\delta-1)\bigr)=\alpha_2m^2+\alpha_1m+\alpha_0\); alors il résulte en comparant des polynômes que \(2\alpha_2=\langle c_1(L)^2\rangle\), que \(\alpha_0=\sum_{i=0}^{2}e_i\binom{\delta-1+i}{i}\) et que \(\alpha_2+\alpha_1+\alpha_0=\sum_{i=0}^{2}a_i\binom{\delta-1+i}{i}\). Soit \(\chi\bigl(L^{-1}(\delta)(n)\bigr)=\sum_{i=0}^{2}\beta_i\binom{n+i}{i}\); alors \(\beta_0=\alpha_2-\alpha_1+\alpha_0\) et par 2.3, \(\beta_1=d(L^{-1}(\delta))+e_1\), ce qui égale \(\beta\). Donc \(a_0=2\alpha_2+2\alpha_0-\sum_{i=1}^{2}a_i\binom{\delta-1+i}{i}-\beta_0\) et comme \(d(L^{-1}(\delta))\le0\), on a \(\beta_0\le\beta(\beta+1)\) par 2.7 (ii); d'où l'assertion.

\textbf{Lemme 2.9.} \emph{On suppose que \(r\ge2\), que \(d(L)\le0\) et que \(X\) soit normal. Alors} :
\begin{enumerate}[label=(\roman*)]
\item \emph{\(h^1(L(-n))=0\) pour \(n\ge2+h_2(a_{r-2},a_{r-1},a_r)\),}
\item \emph{\(h^1(L(n))\le h_r(a_{r-2},a_{r-1},a_r)\) pour tout \(n\),}
\item \emph{\(|a_i|\le g_{r-i}(a_{r-2},a_{r-1},a_r)\) pour \(i=0,\ldots,r\),}
\end{enumerate}
\emph{où \(h_2(X_0,X_1,X_2),h_3(X_0,X_1,X_2),\ldots\) et \(g_0(X_0,X_1,X_2),g_1(X_0,X_1,X_2),\ldots\) sont des polynômes universels à coefficients dans \(\mathbb Q\), qui ne dépendent ni de \(k\) ni de \(X\).}

